problem:  
Let \((M^{n+1}, g)\) be a time‑oriented Lorentzian manifold with smooth, compact, spacelike codimension‑2 surface \(S\subset M\) and future‑directed null geodesic congruence orthogonal to \(S\).  Suppose \((M,g)\) satisfies the **null convergence condition** and that the expansion \(\theta_+(0)=0\) on \(S\).  

Denote by \(J(\lambda)\) the associated **Jacobi tensor map**
\[
J''+\mathcal R(\lambda)J=0,\qquad J(0)=I,\ J'(0)=\chi_0
\]
along each null generator, where \(\chi_0=\chi(0)\) is the initial null shape operator.  The cross‑sectional metric evolves as \(h_\lambda = J^T h_S J\).  

In Riemannian comparison geometry, **Rauch’s theorem** controls \(J^T J\) under curvature inequalities but crucially uses *constant sectional curvature bounds* and *positive‑definite metrics.*  In the Lorentzian null case, \(h_\lambda\) is *degenerate* and its evolution lives in an indefinite bundle over null directions.

**Question:**  
Can one construct a consistent extension of Rauch‑type operator comparison to the **degenerate optical geometry** determined by a null congruence, i.e. establish comparison theorems for \(h_\lambda\) or \(J\) under curvature bounds
\[
\mathcal R_{\min}\le\mathcal R(\lambda)\le\mathcal R_{\max}
\]
that are *stable under the null degeneracy* of \(h_S\)?  
Equivalently, can one define and prove a **Lorentzian Rauch comparison theorem for null directions** giving optimal control of the evolution of degenerate bilinear forms \(h_\lambda\), and identify rigidity cases where equality implies shear‑free and curvature‑constant congruences?  

---

Why is it a "good" problem:  
1. **Profound insight:** The classical Rauch comparison fails directly in null geometries because the metric is degenerate and the standard index‑form arguments collapse.  Developing a Rauch‑type theory in this indefinite, degenerate setting would require re‑conceptualizing basic comparison geometry to include causal structure.  

2. **Bridge between fields:**  The problem bridges **comparison geometry**, **Lorentzian causality**, and **geometric analysis**, requiring analytic reformulation of Jacobi‑field methods in a degenerate bundle and potentially creating a new mathematical framework for null hypersurface analysis.  

3. **Extreme simplicity:**  The question asks: *Can the foundational mechanism of Riemannian comparison survive when the metric degenerates along null directions?*  Its statement is conceptually clear, but its resolution would involve deep structural innovations.  

4. **Potential to open new directions:**  A successful “null Rauch theorem” would inaugurate a **Lorentzian comparison geometry** for causal structures, with implications for singularity formation, horizon rigidity, and causal boundary theory.  Even a proof of impossibility would reveal fundamental obstructions specific to null curvature geometry.  

This version isolates a genuinely unresolved gap—extending comparison theory to degenerate null directions—making it a precise, original, and deeply consequential research problem.